
\assignmenttitle
    {LAB ASSIGNMENT - 03}
    {SCHEDULING ALGORITHMS}

\aim{To implement various scheduling algorithms in C++. (First come first serve, Shortest job first, Shortest running task first, Longest job first, Longest running task first, Non-preemptive priority based, Preemptive priority based, Round robin)}

\appr{Computer System with LINUX/UNIX OS and a C++ compiler installed.}

\sectionlabel{Theory:}
\begin{itemize}
\item\textbf{FCFS (First Come First Serve):} Executes processes in the order of their arrival time. It is non-preemptive, so a process runs until completion once it starts, thus can suffer from "Convoy Effect" where one long process blocks execution of multiple processes.
\item\textbf{SJF (Shortest Job First):} Selects the process with the shortest burst time among the arrived processes. It is non-preemptive, so the selected process runs until completion.
\item\textbf{SRTF (Shortest Remaining Time First):} Selects the process with the shortest remaining CPU time. It is preemptive, so a newly arrived process with a shorter remaining time can interrupt the current process.
\item\textbf{LJF (Longest Job First):} Selects the process with the longest burst time among the arrived processes. It is non-preemptive, so the selected process runs until completion.
\item\textbf{LRTF (Longest Remaining Time First):} Selects the process with the longest remaining CPU time. It is preemptive, so a process with a longer remaining time can interrupt the current process.
\item\textbf{Priority Scheduling (Non-Preemptive):} Selects the highest-priority process among the arrived processes. Once execution starts, the process runs until completion.
\item\textbf{Priority Scheduling (Preemptive):} Selects the process with the highest priority for execution. A newly arrived process with higher priority can interrupt the currently running process.
\item \textbf{Round Robin:} Assigns each process a fixed time quantum and executes them in circular order. If a process does not finish within its quantum, it is placed at the back of the ready queue.
\end{itemize}


\programs

\program{
    Process class and helper functions:
}
\codelabel

\program{
    process.hpp:
}
\inputCodeFile{OS_Assignments/Code/SchedulingAlgorithms/process.hpp}
\program{
    helper.hpp:
}
\inputCodeFile{OS_Assignments/Code/SchedulingAlgorithms/helper.hpp}


\program{
    3.1. First Come First Serve:
}

\codelabel


\inputCodeFile{OS_Assignments/Code/SchedulingAlgorithms/fcfs.cpp}


\outputlabel
\outputimage[0.50\linewidth]{images/OS/scheduling1.png}




\program{
    3.2. Shortest Job First:
}

\codelabel


\inputCodeFile{OS_Assignments/Code/SchedulingAlgorithms/sjf.cpp}


\outputlabel
\outputimage[0.50\linewidth]{images/OS/scheduling2.png}



\program{
    3.3. Longest Running Task First:
}

\codelabel


\inputCodeFile{OS_Assignments/Code/SchedulingAlgorithms/lrtf.cpp}


\outputlabel
\outputimage[0.50\linewidth]{images/OS/scheduling5.png}



\program{
    3.4. Shortest Running Task First:
}

\codelabel


\inputCodeFile{OS_Assignments/Code/SchedulingAlgorithms/srtf.cpp}


\outputlabel
\outputimage[0.50\linewidth]{images/OS/scheduling3.png}



\program{
    3.5. Longest Job First:
}

\codelabel


\inputCodeFile{OS_Assignments/Code/SchedulingAlgorithms/ljf.cpp}


\outputlabel
\outputimage[0.50\linewidth]{images/OS/scheduling4.png}



\program{
    3.6. Priority based non-preemptive:
}

\codelabel


\inputCodeFile{OS_Assignments/Code/SchedulingAlgorithms/non_preemptive_priority.cpp}


\outputlabel
\outputimage[0.50\linewidth]{images/OS/scheduling6.1.png}
\outputimage[0.50\linewidth]{images/OS/scheduling6.2.png}



\program{
    3.7. Priority based preemptive:
}

\codelabel


\inputCodeFile{OS_Assignments/Code/SchedulingAlgorithms/preemptive_priority.cpp}


\outputlabel
\outputimage[0.50\linewidth]{images/OS/scheduling7.1.png}
\outputimage[0.50\linewidth]{images/OS/scheduling7.2.png}

\sectionlabel{Learning Outcomes:}
\begin{itemize}
\item{Working principles of different CPU scheduling algorithms.}
\item{Implementation of preemptive and non-preemptive scheduling techniques.}
\item{Compare scheduling algorithms based on their execution order and average turnaround and waiting times.}
\end{itemize}

\concl{Various CPU scheduling algorithms, including FCFS, SJF, SRTF, LJF, LRTF, Priority Scheduling, and Round Robin were implemented successfully.}
